\documentclass[12pt]{article}

\input{../../_assets/preambulo.tex}


\begin{document}

    % 1. Foto de fondo
    % 2. Título
    % 3. Encabezado Izquierdo
    % 4. Color de fondo
    % 5. Coord x del titulo
    % 6. Coord y del titulo
    % 7. Fecha

    
    \input{../../_assets/portada}
    \portadaExamen{ffccA4.jpg}{Álgebra I\\Examen XIII}{Álgebra I. Examen XIII}{MidnightBlue}{-8}{28}{2025}{David Sánchez Muñoz}

    \begin{description}
        \item[Asignatura] Álgebra I.
        \item[Curso Académico] 2024-2025.
        \item[Grado] Doble Grado Informática-Matemáticas.
        \item[Grupo] Único.
        \item[Profesor] María Pilar Carrasco Carrasco.
        \item[Descripción] Examen Ordinario.
        \item[Fecha] 21 de enero de 2025.
        \item[Duración] 3 horas.
    
    \end{description}
    \newpage


    % ------------------------------------
    
    \begin{ejercicio}[$2.5$ puntos]
        Demostrar:
        \begin{enumerate}[label=\arabic*.]
            \item ($1.25$ puntos) Todo DI en el que todo elemento no nulo ni unidad es producto de irreducibles y todo irreducible es primo, es un DFU.
            \item ($1.25$ puntos) Sea $R$ una relación de equivalencia definida en un conjunto $S$, entonces para cualesquiera $a, b \in S$,
            \[
            aRb \sii [a] = [b].
            \]
        \end{enumerate}
    \end{ejercicio}

    \begin{ejercicio}[$2.5$ puntos]
        Dado el sistema de congruencias en $\bb{Z}[\sqrt{-2}]$
        \[
        \begin{cases}
            x \equiv 1+2\sqrt{-2} & \pmod{2-3\sqrt{-2}} \\
            x \equiv 3 & \pmod{1+\sqrt{-2}}
        \end{cases}
        \]
        \begin{enumerate}
            \item Demostrar que tiene solución sin resolverlo.
        
            \item Resolverlo, dando una solución óptima y la solución general.
        
            \item ¿Es posible encontrar una solución cuya parte real sea $a = 3$? En caso afirmativo, describir todas ellas.
        \end{enumerate}
    \end{ejercicio}

    \begin{ejercicio}[$2.5$ puntos]
        Factorizar en el anillo $\bb{Z}[x]$, y en $\bb{Q}[x]$ como producto de irreducibles mónicos, los dos siguientes polinomios:
        \begin{enumerate}[label=\arabic*.]
            \item ($1.25$ puntos) $6x^7 + 30x^6 + 6x^2 + 36x + 30$.
            \item ($1.25$ puntos) $3x^6 - 4x^5 + 49x^4 - 64x^3 + 16x^2 + 3x - 1$.
        \end{enumerate}
    \end{ejercicio}

    \begin{ejercicio}[$2.5$ puntos]~
        \begin{enumerate}[label=\arabic*.]
            \item ($0.75$ puntos) En el anillo $\bb{Z}_7[x]$ calcular el resto de dividir $x^{1321} + 5$ entre el polinomio $x + 3$.
            \item ($1.75$ puntos) Sea $\alpha = 1 + \sqrt{2} \in \bb{R}$ y sea
            \[
            I = \{ f \in \bb{Q}[x] \mid f(\alpha) = 0 \}
            \]
            Demostrar que $I$ es un ideal no nulo de $\bb{Q}[x]$.
            
            Encontrar un polinomio $h \in \bb{Q}[x]$ tal que $I = h\bb{Q}[x]$, el ideal principal generado por $h$.
        \end{enumerate}
    \end{ejercicio}

    % \newpage
    % \setcounter{ejercicio}{0} % Reiniciar contador de ejercicios
    % \noindent
    % \textbf{Solución.}

\end{document}